\newpage
\section{Lecture-11}
\subsection{Inverse Function Theorem}
The phrase \emph{inverse function} appears first in elementary mathematics in a set-theoretic form: a function
\[
f:X\to Y
\]
has an inverse if and only if it is bijective. This is correct, but in analysis and geometry it is not the whole story. For maps between Euclidean spaces,
\[
f:\R^m\to\R^n,
\]
we usually want more than merely an inverse as a rule of correspondence between sets. We want the inverse to preserve the analytic structure of the problem: continuity, differentiability, smoothness, and the local geometry of the spaces involved.

The inverse function theorem explains when a smooth map behaves \emph{locally} like a reversible change of coordinates.
Let \(f:X\to Y\) be a function between sets.

\begin{proposition}
The following are equivalent:
\begin{enumerate}[label=(\arabic*)]
    \item \(f\) has an inverse function \(f^{-1}:Y\to X\),
    \item \(f\) is bijective,
    \item for each \(y\in Y\), there exists a unique \(x\in X\) such that \(f(x)=y\).
\end{enumerate}
\end{proposition}
This statement is purely about matching elements. It does not know anything about: continuity, open sets, differentiability, tangent vectors, local coordinates.
So from the point of view of analysis, the set-theoretic inverse is too weak. A bijection can be very irregular.
For maps between spaces carrying more structure, there are different notions of invertibility:
\begin{itemize}
    \item \textbf{Set-theoretic inverse:} \(f\) is bijective.
    \item \textbf{Topological inverse:} \(f\) is a homeomorphism.
    \item \textbf{Smooth inverse:} \(f\) is a diffeomorphism.
\end{itemize}
The inverse function theorem belongs to the smooth category. Suppose \(f:\R^n\to\R^n\) is smooth. Even if \(f\) is globally bijective, that still does not guarantee that \(f^{-1}\) is smooth.
\begin{example}
Consider
\[
f:\R\to\R,\qquad f(x)=x^3.
\]
This map is bijective and smooth. Its inverse is
\[
f^{-1}(y)=\sqrt[3]{y}.
\]
The inverse is continuous, but it is not differentiable at \(y=0\). Indeed,
\[
\frac{d}{dy}\bigl(y^{1/3}\bigr)=\frac{1}{3y^{2/3}},
\]
which is not defined at \(0\).
\end{example}
So even in the easiest case, global bijectivity is not enough for a smooth inverse. Something more must be required.
The inverse function theorem is not a theorem about global invertibility. It is a theorem about \emph{local} invertibility.
\begin{example}
Let
\[
f:\R\to\R,\qquad f(x)=x^3-x.
\]
This function is not globally injective. However,
\[
f'(x)=3x^2-1.
\]
At any point where \(3x^2-1\neq 0\), the function is locally invertible.
\end{example}
Thus local invertibility is controlled by the derivative, not by global behavior.
~\\
Let \(f:U\to\R^n\) be differentiable, where \(U\subseteq \R^n\) is open, and let \(a\in U\). Then for small \(h\),
\[
f(a+h)=f(a)+Df(a)\,h+\text{error},
\]
where the error is small compared with \(\|h\|\). So near \(a\), the function behaves like its derivative \(Df(a)\), which is a linear map
\[
Df(a):\R^n\to\R^n.
\]
This leads to the fundamental idea:
\begin{center}
\emph{To understand whether \(f\) is locally invertible near \(a\), first ask whether the linear map \(Df(a)\) is invertible.}
\end{center}

If \(Df(a)\) is invertible, then its action on small balls is easy to visualize: it may stretch, rotate, reflect, or shear, but it does not collapse any direction.

If \(Df(a)\) is singular, then some infinitesimal direction is flattened, and local invertibility may fail.
\begin{theorem}[Inverse Function Theorem]
Let \(U\subseteq \R^n\) be open, and let
\[
f:U\to\R^n
\]
be a \(C^1\) map. Suppose that for some \(a\in U\), the derivative \(Df(a)\) is invertible. Equivalently,
\[
\det Df(a)\neq 0.
\]
Then there exist open neighborhoods \(U_0\subseteq U\) of \(a\) and \(V_0\subseteq \R^n\) of \(f(a)\) such that
\[
f:U_0\to V_0
\]
is a bijection, and its inverse
\[
f^{-1}:V_0\to U_0
\]
is also \(C^1\).

Moreover,
\[
D(f^{-1})(f(a))=(Df(a))^{-1}.
\]
\end{theorem}
The theorem says that invertibility of the derivative at one point forces invertibility of the nonlinear map on small neighborhoods.
\section*{Exercises}
\begin{example}[A linear map]
Let
\[
f(x,y)=(2x+y,\;x-y).
\]
Then
\[
Df(x,y)=
\begin{pmatrix}
2 & 1\\
1 & -1
\end{pmatrix},
\qquad
\det Df(x,y)=-3\neq 0.
\]
Hence \(f\) is locally invertible everywhere. Since this is a linear map with an invertible matrix, it is in fact globally invertible.
\end{example}

\begin{example}[Polar coordinates away from the origin]
Consider
\[
f(r,\theta)=(r\cos\theta,\;r\sin\theta),
\qquad r>0.
\]
Its Jacobian matrix is
\[
Df(r,\theta)=
\begin{pmatrix}
\cos\theta & -r\sin\theta\\
\sin\theta & r\cos\theta
\end{pmatrix},
\]
so
\[
\det Df(r,\theta)=r.
\]
For \(r>0\), this is nonzero. Hence \(f\) is locally invertible wherever \(r>0\).

However, \(f\) is not globally injective, because \(\theta\) and \(\theta+2\pi\) determine the same point.
\end{example}

\begin{example}[A smooth bijection with non-smooth inverse]
Let
\[
f(x)=x^3.
\]
Then \(f\) is bijective and smooth, but
\[
f'(0)=0.
\]
The inverse is not differentiable at \(0\). So \(f\) is not a local diffeomorphism at the origin.
\end{example}

\begin{example}[Failure when dimensions do not match]
Define
\[
f:\R^2\to\R,\qquad f(x,y)=x^2+y^2.
\]
Then
\[
Df(x,y)=\begin{pmatrix}2x & 2y\end{pmatrix},
\]
which is a \(1\times 2\) matrix. There is no square Jacobian here, hence no usual local inverse between open subsets of Euclidean spaces.
\end{example}
\subsection{Implicit Function Theorem}
Let's look $x^2+y^2=1$ closely. Is this a graph of a function? It's not.
No, it's not. It's a relation between $x$ and $y$. But we want to study function so far. 
At a neighborhood of $(x,y)$ except $(0,1)$ and $(0,-1)$, we can solve for $y$ as a function of $x$. 
$y=\sqrt{1-x^2}$ or $y=-\sqrt{1-x^2}$ are explicit functions that give the same relation between $x$ and $y$
as $x^2+y^2=1$. This is true locally near any point except $(0,1)$ and $(0,-1)$.

The implicit function theorem provides conditions under which a relationship of the from $F(x,y)=0$ can be rewritten as
a function $y=f(x)$ locally. 
\begin{example}
    Consider $y^5+y^3+y+x=0$. It is trivial to write $x=f(y)$, but can we write $y=f(x)$?
    But there is no explicit formula for $y$ in terms of $x$, because there is no formula for the roots of a general quintic polynomial.
    In general, there is no formula for finding roots of polynomials of degree 5 or higher (Abel, Galois). 
    Fix $x_0$, define $\varphi(y)=y^5+y^3+y+x_0$. Then,
\[
\varphi'(y)=5y^4+3y^2+1>0,\qquad\forall y\in\R.
\] 
Which implies $\varphi$ is strictly increasing. In Calculus-I, we know that any odd degree polynomial has at least one real root.
But since $\varphi$ is strictly increasing, it has exactly one real root. So we can write $y=f(x)$ locally near any point $(x_0,y_0)$, give me 
$x_0$, we can find a unique $y_0$ such that $y_0^5+y_0^3+y_0+x_0=0$. We found a function $y=f(x)$, but $f$ is implicit.       
\end{example} 
\begin{theorem}
Denote $\vec x=(x_1,\ldots,x_n)$. Let $F(\vec x,y)\in C^1$ in a neighborhood of $N_0(\vec x_0,y_0)$ such that
\begin{enumerate}
    \item $F(x_0,y_0)=0$,
    \item $\dfrac{\partial F}{\partial y}(x_0,y_0)\neq 0$.
\end{enumerate}
Then there exists a neighborhood of \((x_0,y_0)\) in which there is a unique (implicit) function
$y=f(\vec x)$ such that
\begin{enumerate}
    \item[(i)] \(y_0=f(\vec x_0)\),
    \item[(ii)] \(F\bigl(\vec x,f(\vec x)\bigr)=0\) for every \(x\) in the neighborhood,
    \item[(iii)] And the derivative of \(f\) is given by,
    \[
    \frac{\partial f}{\partial x_i}=-\frac{\dfrac{\partial F}{\partial x_i}}{\dfrac{\partial F}{\partial y}}
    \]
    in the neighborhood of $(\vec x_0,y_0)$.
\end{enumerate}
\end{theorem} 
\begin{example}
Given that $3x^2y - yz^2 - 4xz - 7 = 0$. We will show that near \((-1,1,2)\) we can write
$y = f(x,z)$, and we will find $\frac{\partial f}{\partial x}(-1,2)$.
Here,
\[
F(x,y,z) = 3x^2y - yz^2 - 4xz - 7 \in C^1
\]
And $F(-1,1,2) = 0$. Moreover,
\[
\frac{\partial F}{\partial y} = 3x^2 - z^2,
\qquad
\frac{\partial F}{\partial y}(-1,1,2) = -1 \ne 0.
\]
Thus, by the implicit function theorem, there exists a neighborhood of \((-1,1,2)\) in which there is a unique function \(y=f(x,z)\in C^1\) such
that \(f(-1,2)=1\) and \(F\bigl(x,f(x,z),z\bigr)=0\) for every \((x,z)\) in the neighborhood. Moreover,
\[
\frac{\partial f}{\partial x}(-1,2) = -\frac{\frac{\partial F}{\partial x}}{\frac{\partial F}{\partial y}}(-1,1,2) = -\frac{6x^2y - 4z}{3x^2 - z^2}(-1,1,2) = -\frac{6(-1)^2(1) - 4(2)}{3(-1)^2 - (2)^2} = -\frac{6 - 8}{3 - 4} = -\frac{-2}{-1} = -2.
\]
In fact, 
\begin{align*}
    3x^2y-yz^2&= 4xz+7\\
    y(3x^2-z^2) &= 4xz+7\\
    y&=\frac{4xz+7}{3x^2-z^2}.
\end{align*}
In this example, we can explicitly solve for \(y\) in terms of \(x\) and \(z\), but the implicit function theorem guarantees the existence of such a function even when an explicit formula is not available. The key point is that the condition \(\frac{\partial F}{\partial y}(-1,1,2) \neq 0\) ensures that we can locally solve for \(y\) as a function of \(x\) and \(z\).
In this example, we can write $z=f(x,y)$ as well, but the theorem fails at \((-1,1,2)\) because
\[\frac{\partial F}{\partial z}(-1,1,2) = 0.\]
\end{example}
\begin{example}
\[
\text{ex.}\qquad F(x,y)=(x-y)^3
\]
\[
F(x,y)=0 \iff y=x
\]
There is an explicit function at any point.
However, at \((0,0)\),
\[
\frac{\partial F}{\partial x}(0,0)=3(x-y)^2\big|_{(0,0)}=0,
\]
\[
\frac{\partial F}{\partial y}(0,0)=-3(x-y)^2\big|_{(0,0)}=0.
\]
The theorem doesn't apply.
\end{example}
Let's prove the implicit function theorem for the case of two variables only.  
\begin{proof}
    $\frac{\partial F}{\partial y}(x_0,y_0)\neq 0$ implies w.l.o.g. $\frac{\partial F}{\partial y}(x_0,y_0)>0$.
    Since, $F\in C^1$, $\frac{\partial F}{\partial y}$ is continuous, so $\frac{\partial F}{\partial y}>0$ 
    at a neighborhood of $(x_0,y_0)$. What do we know about $F(x_0,y)$ function? Here we fixing $x_0$,
    and view $F$ as a function of $y$. We know that $\frac{\partial F}{\partial y}>0$ at a neighborhood of $(x_0,y_0)$, which implies
    $F(x_0,y)$ is a strictly increasing function. $\exists y_1$ sucht $F(x_0,y_1)>0$ and $\exists y_2$ such that $F(x_0,y_2)<0$.
    Since $F\in C^1$, $F$ is continuous. Which implies $\forall x$ near $x_0$ (releasing our first variable),
    \[ F(x,y_1)>0, \qquad F(x,y_2)<0\]
    For such an $x$ near $x_0$, since $\frac{\partial F}{\partial y}>0$, $(x,y)$ is strictly increasing as a function of $y$, therefore $! \exists y$ such that $F(x,y)=0$.
    This proves that $y=f(x)$ exists and is unique. Now,
    \begin{align*}
        F(x,f(x))&=0\\
        \frac{\partial F}{\partial x} \cdot 1+ \frac{\partial F}{\partial y} \frac{d f}{d x} &= 0
        \implies \frac{d f}{d x} = -\frac{\frac{\partial F}{\partial x}}{\frac{\partial F}{\partial y}}.
    \end{align*}
\end{proof}
\begin{theorem}
    $$
    \begin{cases}
        F_1(x_1,\cdots,x_n,z_1,\cdots,z_m)&=0\\
        \vdots\\
        F_m(x_1,\cdots,x_n,z_1,\cdots,z_m)&=0
    \end{cases}
    $$
    where $F_i\in C^1$ for $i=1,\cdots,m$. Let 
\[\triangle=\det\begin{pmatrix}
\frac{\partial F_1}{\partial z_1} & \cdots & \frac{\partial F_1}{\partial z_m}\\
\vdots & \ddots & \vdots\\
\frac{\partial F_m}{\partial z_1} & \cdots & \frac{\partial F_m}{\partial z_m}
\end{pmatrix}\]
If $\triangle\neq 0$ at a point $(\vec x^0,\vec z^0)$, then near that point, there are unique functions $z_i = f_i(x_1,\cdots,x_n)$ for $i=1,\cdots,m$ such that the system of equations is satisfied.
and the partial derivatives of $f_i$'s can be computed using implicit differentiation.
\end{theorem}
\begin{example}
Consider the system of equations:
\[
\begin{cases}
    xu+yvu^2&=2\\
    xu^3+y^2v^4&=2
\end{cases}
\]
Let's show that near $(x,y,u,v)=(1,1,1,1)$, we can express $u$ and $v$ as functions of $x$ and $y$.
\[
\begin{cases}
F_1(x,y,u,v) &= xu+yvu^2-2\\
F_2(x,y,u,v) &= xu^3+y^2v^4-2
\end{cases}
\]
Here, \[
\triangle = \det\begin{pmatrix}
\frac{\partial F_1}{\partial u} & \frac{\partial F_1}{\partial v}\\
\frac{\partial F_2}{\partial u} & \frac{\partial F_2}{\partial v}
\end{pmatrix} = \det\begin{pmatrix}x+2yvu & yu^2\\
3xu^2 & 4y^2v^3 \end{pmatrix}.
\]
Evaluating at the point $(1,1,1,1)$ gives:
\[\triangle(1,1,1,1) = \det\begin{pmatrix}1+2(1)(1) & 1(1)^2\\
3(1)(1)^2 & 4(1)^2(1)^3 \end{pmatrix} = \det\begin{pmatrix}3 & 1\\
3 & 4 \end{pmatrix} = 3\cdot4 - 1\cdot3 = 12 - 3 = 9 \neq 0.\]
By the implicit function theorem, there exist unique functions \(u=f_1(x,y)\) and \(v=f_2(x,y)\) near the point $(1,1,1,1)$.
Let's find $\frac{\partial u}{\partial x}(1,1)$:
\[
\begin{cases}
    u+x\frac{\partial u}{\partial x}+y\frac{\partial v}{\partial x}u^2+2yvu\frac{\partial u}{\partial x}&=0\\
    u^3+3xu^2\frac{\partial u}{\partial x}+y^2\cdot 4v^3\frac{\partial v}{\partial x}&=0
\end{cases} \implies 
\begin{cases}
    3\frac{\partial u}{\partial x}+\frac{\partial v}{\partial x}&=-1\\
    3\frac{\partial u}{\partial x}+4\frac{\partial v}{\partial x}&=-1
\end{cases}
\]  
\end{example}